\section{Challenge Input Data}
\label{sec:challenge_input_data}

The preparation of the challenge data is described in the appendices.
The data are available as a \texttt{tar} archive that is downloaded and
unpacked as part of the \texttt{pz\_data\_challenge} setup procedure.

Each task set in the data challenge has an associated set of files.
Typically these will be a collection of training files that contain
photometric data and reference redshifts, and a second set of files
that contain photometric data but do not include redshifts. Each
task set will involve estimating something about the redshifts or
redshift distributions in the test files.

Typically there will be several training and test files for a
particular task set, covering different scenarios and using different
input simulations.


\subsection{Input data format}

The input data for the challenge are presented in HDF5 files.
The naming convention for the files is
\texttt{\{challenge\}\_\{taskset\}\_\{simulation\}\_\{label\}\_\{scenario\}.hdf5}.
The meanings of the various \allowbreak fields are described in
Tab.~\ref{tab:file_fields}. The columns in the files are described
in Tab.~\ref{tab:columns}.  We note that we use \texttt{np.nan} to in
the magnitdude columns to signify non-detections.

\begin{table}[ht]
\centering
\caption{\label{tab:file_fields}\small Fields in the input file names.}
\begin{tabular}{ll}
  \hline\hline  
  Field & Description \\ \hline
  challenge & Challenge associated with file (``pz\_challenge'') \\
  taskset & Task set associated with file (e.g., ``taskset\_1'') \\
  simulation & Simulation used to produce file (``cardinal'' or ``flagship'') \\
  label & File label (e.g., ``test'', ``training'') \\
  scenario & Data scenario (e.g., ``1yr'', ``10yr'') \\
\hline\hline
\end{tabular}
\end{table}


\begin{table}[ht]
\centering
\caption{\label{tab:columns}\small Contents of input files.}
\begin{tabular}{ll}
  \hline\hline
  Column & Description \\ \hline  
  redshift & True redshift (training files only) \\
  ra & Right ascension (training files only) \\
  dec & Declination (training files only) \\
  object\_id & Unique object ID\\
  mag\_\{band\}\_lsst & Magnitude in LSST \{band\} \\
  mag\_\{band\}\_lsst\_err & Magnitude uncertainty in LSST \{band\} \\
  mag\_\{band\}\_roman & Magnitude in Roman \{band\} \\
  mag\_\{band\}\_roman\_err & Magnitude uncertainty in Roman \{band\} \\
\hline\hline
\end{tabular}
\end{table}

We note that the \texttt{table-io} package~
\cite{tables_io} installed with
\texttt{pz\_data\_challenge} provided a command line interface to
convert files from \texttt{hdf5} format to other formats such as
\texttt{parquet} tables or \texttt{pandas} data frames.

\begin{verbatim}
# convert a hdf5 file to pandas dataframe in a parquet file
tables-io convert
  --input public/pz_challenge_taskset_1_cardinal_test_10yr.hdf5
  --output public/pz_challenge_taskset_1_cardinal_test_10yr.pq
\end{verbatim}

